%=====================================================================
% Dokument: EG-1 – Kugel-Schatten-Prinzip (MKO-konform)
% Version: 1.0   (Juni 2026)
%=====================================================================

\documentclass[11pt,a4paper]{article}

% ---------- Grundpakete ----------
\usepackage[utf8]{inputenc}
\usepackage[T1]{fontenc}
\usepackage[german]{babel}
\usepackage{amsmath,amssymb,amsthm}
\usepackage{geometry}
\usepackage{parskip}
\usepackage{enumitem}
\usepackage{booktabs}
\usepackage{xcolor}
\usepackage{underscore} % FIX: Erlaubt Unterstriche im Dateinamen (Jobname) ohne Absturz
\usepackage{hyperref}
\usepackage{csquotes}   % schöne Anführungszeichen

% ---------- Layout ----------
\geometry{margin=2.5cm}
\setlist{itemsep=0.2em, topsep=0.2em}
\hypersetup{
    colorlinks=true,
    linkcolor=blue,
    urlcolor=blue,
    pdftitle={EG-1 – Kugel-Schatten-Prinzip (MKO)},
    pdfauthor={Kritischer Dialog}
}

% ---------- Theorem-Umgebungen ----------
\theoremstyle{definition}
\newtheorem{konvention}{Konvention}[section]
\newtheorem{remark}{Bemerkung}[section]
\newtheorem{methodennote}{Methodennote}[section]
\newtheorem{definition}{Definition}[section]

\theoremstyle{plain}
\newtheorem{axiom}{Axiom}[section]
\newtheorem{proposition}{Proposition}[section]
\newtheorem{hypothesis}{Hypothese}[section]

% ---------- Titel ----------
\title{
  \textbf{Kugel-Schatten-Prinzip}\\[0.4em]
  {\large Formalisierung der Analogie im Rahmen der MKO}\\[0.5em]
  {\normalsize Version 1.0}
}
\author{
  \textbf{Kritischer Dialog:}\\
  DeepSeek (Seki), Apeiron, Grok, Gemini, Claude, ChatGPT, Manus\\
  Gemeinsames Logbuch eines offenen Forschungsprogramms
}
\date{Juni 2026}

\begin{document}
\maketitle

\begin{abstract}
Wir entwickeln die \textbf{Kugel-Schatten-Analogie} – das zentrale Bild des
Ontophänomenalismus (OP) – zu einer formalen mathematischen Struktur.
Dabei wird die phänomenale Grundstruktur $B$ als „Kugel“ $\mathcal{K}$ und
die physikalische Manifestation $A$ als „Schatten“ $\mathcal{S}$ dargestellt.
Durch die Einführung eines \textbf{Projektionsoperators} $P\colon\mathcal{K}\to\mathcal{S}$
zeigen wir, dass das OP-Framework die fünf-Ebenen-Sprachregel des
\textbf{Methodologischen Konventionen-Frameworks (MKO)} einhält und die
Zirkel-Kontrolle transparent macht.
\end{abstract}

\tableofcontents
\newpage

%=====================================================================
\section{Einleitung}
%=====================================================================
Das \emph{Hard Problem of Consciousness} wird im OP durch das Bild
der \textbf{Kugel-Schatten-Relation} erklärt: das phänomenale Erleben ist die
Kugel, die physikalische Welt ihr Schatten. Im Folgenden wird diese
Analogie in mathematische Formeln überführt, wobei jede Aussage nach den
MKO-Regeln (5 Ebenen, Kennzeichnungen) gekennzeichnet wird.

%=====================================================================
\section{Grundlegende Begriffe}
%=====================================================================

\subsection{Phänomenale Grundstruktur (Kugel)}
\begin{definition}[Kugel]
\textbf{Kugel $\mathcal{K}$}: Der ontologische Grundbegriff des OP,
der das gesamte, unveränderliche Feld des bewussten Erlebens
(z.\,B. das „Ur-Bewusstsein“) bezeichnet. Sie ist ein topologisch
kompaktes, nicht-diskretes Hilbert-ähnliches Objekt mit innerer
Selbst-Referenz.
\end{definition}

\begin{remark}
Ebene 3 – Definition.
\end{remark}

\subsection{Physikalische Manifestation (Schatten)}
\begin{definition}[Schatten]
\textbf{Schatten $\mathcal{S}$}: Das physikalische, messbare
System (z.\,B. das Gehirn, neuronale Netze, makroskopische Felder),
das aus der Kugel hervorgeht. Es ist ein Banach-Raum, dessen Elemente
durch experimentelle Messgrößen (EEG-Signale, fMRT, etc.) repräsentiert
werden können.
\end{definition}

\begin{remark}
Ebene 3 – Definition.
\end{remark}

\subsection{Projektionsoperator}
\begin{definition}[Projektionsoperator]
\textbf{Projektionsoperator $P$}: Eine lineare Abbildung
\[
   P\colon\mathcal{K}\longrightarrow\mathcal{S}
\]
mit den Eigenschaften
\[
   P^{2}=P\qquad\text{(Idempotenz)}\tag{1}
\]
\[
   \ker P\neq\{0\}\qquad\text{(nicht-injektiv)}\tag{2}
\]
\end{definition}

\begin{remark}
Ebene 3 – Definition.
\end{remark}

%=====================================================================
\section{Mathematische Ausarbeitung der Analogie}
%=====================================================================

\subsection{Geometrische Intuition}
\begin{methodennote}[A1]
Stellen wir uns $\mathcal{K}$ als eine dreidimensionale Kugel in
$\mathbb{R}^{3}$ vor. Die Projektion $P$ ist das orthogonale
Abbilden auf eine zweidimensionale Ebene $\mathcal{S}$ (der „Boden“). Jeder
Punkt $x\in\mathcal{K}$ besitzt ein Bild $Px\in\mathcal{S}$, das wir als
\textbf{Schatten} bezeichnen. Da $P$ nicht-injektiv ist, besitzen
unterschiedliche Punkte von $\mathcal{K}$ denselben Schatten – das erklärt,
warum unterschiedliche innerweltliche Zustände (z.\,B. unterschiedliche Qualia)
auf das gleiche physikalische Messsignal reduzieren können.
\end{methodennote}

\begin{remark}
Ebene 4 – Analogie.
\end{remark}

\subsection{Formale Eigenschaften}
\begin{itemize}
  \item \textbf{Idempotenz (1)}: $P$ ist ein Projektor; eine wiederholte
        Anwendung ändert das Ergebnis nicht mehr. In physikalischen
        Begriffen bedeutet das, dass das physikalische System ein
        \emph{stabiler} Ausdruck der phänomenalen Struktur ist.
  \item \textbf{Nicht-Injektivität (2)}: $\dim\ker P>0$; es gibt
        nicht-triviale Subräume von $\mathcal{K}$, die denselben Schatten
        erzeugen. Das ist das mathematische Pendant zur \textbf{Mensch-Schatten-Mehrdeutigkeit}
        (verschiedene innere Zustände $\rightarrow$ identische neuronale Aktivität).
  \item \textbf{Bild von $P$}: $\operatorname{im}P=\mathcal{S}$ ist ein
        abgeschlossener Unterraum von $\mathcal{S}$. Damit ist jede
        physikalische Messgröße ein \emph{Projektionsresultat} aus
        $\mathcal{K}$.
\end{itemize}

\begin{remark}
Ebene 2 – Bemerkung zu den formalen Eigenschaften.
\end{remark}

\subsection{Kausal-logische Relation}
\begin{axiom}[R1]
\textbf{Kausalität umgekehrt}: Aus $P(x)=y$ folgt \emph{logisch},
dass $y$ ein Schatten von $x$ ist, nicht umgekehrt. Formal:
\[
   \forall x\in\mathcal{K}\;\exists y\in\mathcal{S}\;(y=P(x))
   \;\Longrightarrow\;
   \forall y\in\mathcal{S}\;\exists x\in\mathcal{K}\;(y=P(x)).
\]
\end{axiom}

\begin{remark}
Ebene 1 – Rekonstruktives Axiom (R).
\end{remark}

\subsection{Postulat der ontologischen Grundentscheidung}
\begin{axiom}[P1]
\textbf{Ontologisches Postulat}: Das \textbf{Kugel-Schatten-Prinzip}
sei die \emph{grundlegende} Relation zwischen phänomenaler
Existenz und physikalischer Manifestation. Jede alternative
Theorie, die $A$ ohne $B$ zulässt, ist eine metaphysische Zusatzannahme.
\end{axiom}

\begin{remark}
Ebene 3 – Postulat (P).
\end{remark}

\subsection{Hypothese zur Empirischen Testbarkeit}
\begin{hypothesis}[H1]
\textbf{Empirische Hypothese}: Für jedes messbare physikalische
System $y\in\mathcal{S}$ gibt es mindestens ein nicht-trivales
phänomenales Element $x\in\mathcal{K}$ mit $P(x)=y$; umgekehrt lässt
sich jedes $x\in\mathcal{K}$ eindeutig einem $y\in\mathcal{S}$ zuordnen.
\end{hypothesis}

\begin{remark}
Ebene 5 – Hypothese (testbar).
\end{remark}

\subsection{Methodennotiz zu Gegenargumenten}
\begin{methodennote}[G1]
\textbf{Mögliche Gegenargumente}:
\begin{enumerate}[label=(\alph*)]
  \item \textbf{Korrelation $\neq$ Kausalität}: Die Projektion $P$ ist
        nur ein \emph{mathematisches Abbilden}, kein kausaler
        Mechanismus.
  \item \textbf{Verborgene physikalische Mechanismen}: Es könnte
        Systeme geben, die $A$ besitzen, ohne dass ein $x\in\mathcal{K}$
        existiert.
\end{enumerate}
\end{methodennote}

\begin{remark}
Ebene 2 – Methodennotiz.
\end{remark}

\subsection{Reaktion auf Gegenargumenten (Proposition)}
\begin{proposition}[P2]
\textbf{Reaktion}: Jede alternative Theorie muss ein
\emph{Projektionsschema} $P'$ besitzen, das die Eigenschaften (1)–(2)
erfüllt und die empirische Hypothese $H1$ reproduziert. Fehlt ein
solches $P'$, bleibt die Theorie eine metaphysische Zusatzannahme.
\end{proposition}

\begin{remark}
Ebene 1 – Proposition (logische Konsequenz).
\end{remark}

%=====================================================================
\section{Zusammenfassung}
%=====================================================================
\begin{enumerate}[label=(\arabic*)]
  \item \textbf{Kugel $\mathcal{K}$} und \textbf{Schatten $\mathcal{S}$}
        sind formal definierte Strukturen (Definitionen K1, S1).
  \item \textbf{Projektionsoperator $P$} verbindet beide Räume und
        erfüllt Idempotenz und Nicht-Injektivität (Definition P1, Formeln 1–2).
  \item \textbf{Rekonstruktives Axiom R1} stellt die logische Richtung
        $B\to A$ (Kugel $\rightarrow$ Schatten) klar.
  \item \textbf{Postulat P1} wählt diese Relation als ontologischen
        Ausgangspunkt des OP.
  \item \textbf{Hypothese H1} macht die Aussage empirisch prüfbar.
  \item \textbf{Methodennotiz G1} und \textbf{Proposition P2} zeigen,
        wie mögliche Gegenargumente in das MKO-Framework eingebettet
        und durch die Zirkel-Kontrolle adressiert werden.
\end{enumerate}

Damit ist das \textbf{Kugel-Schatten-Prinzip} nicht mehr nur ein Bild,
sondern ein formal definiertes, testbares Gerüst, das die
\textbf{Methodologischen Konventionen (MKO)} vollständig respektiert.

\bigskip
\end{document}

